% =========================================================================
% Capítulo 10 — Descripción de tareas, fases, equipos o procedimientos y
% diagrama de Gantt/cronograma (si procede)
% Metodología de trabajo (decisión SCRUM vs. tradicional/RUP, justificada)
% y las tareas/fases del desarrollo con su cronograma.
% OJO: el Gantt del TFG en sí (obligatorio) tiene su propio capítulo (13).
% =========================================================================

\chapter{Descripción de tareas, fases y procedimientos y diagrama de Gantt}
\label{ch:tareas}

Este capítulo justifica la metodología con la que se planifica y ejecuta el trabajo, define cómo se estima el esfuerzo y presenta el \eng{Product Backlog} inicial y las fases del desarrollo con su cronograma.
La representación gráfica del cronograma completo del \ac{TFG} (incluida la memoria y los hitos académicos) se recoge en el capítulo \ref{ch:gantt-tfg}.

\section{Metodología de trabajo}
\label{sec:metodologia}

\subsection{Restricciones que la metodología debe satisfacer}
\label{sec:metodologia-restricciones}

La elección no parte de preferencias abstractas, sino de las condiciones reales de este proyecto:

\begin{itemize}
  \item \textbf{Alcance emergente.}
        El modelo de datos y las entidades se diseñan de forma incremental y meditada durante el desarrollo, una decisión ya tomada del proyecto; una red completa de tareas definida por adelantado quedaría obsoleta en semanas.
  \item \textbf{Disponibilidad variable y conocida.}
        La dedicación de septiembre será sensiblemente menor por motivos personales (riesgo R2, capítulo \ref{ch:riesgos}); la metodología debe absorber esa variación sin romperse.
  \item \textbf{Una única persona.}
        No hay roles que repartir ni reuniones que sincronizar; toda ceremonia debe justificar su coste ante un equipo de uno.
  \item \textbf{Dos entregables simultáneos.}
        Producto y memoria evolucionan a la vez, así que la redacción debe ser trabajo planificado de primera clase, no un resto.
  \item \textbf{Telemetría para esta memoria.}
        El método debe producir por sí mismo los datos de seguimiento (horas, velocidad, desviaciones) que alimentan este capítulo y el \ref{ch:gantt-tfg}.
  \item \textbf{Marco académico.}
        Unas 360 horas de dedicación de referencia (12 créditos \ac{ECTS}) en un periodo de ejecución objetivo de unos tres meses (capítulo \ref{ch:objetivos}).
\end{itemize}

\subsection{Alternativas consideradas y descartadas}
\label{sec:metodologia-alternativas}

\paragraph{Metodologías predictivas (cascada, \acs{RUP}) con planificación por precedencias.}
El enfoque tradicional visto en la asignatura de Gestión de Proyectos (ciclos tipo \ac{RUP} y técnicas de planificación temporal como \ac{PERT}, Roy o los diagramas de precedencias) exige descomponer el proyecto completo en tareas con duraciones y dependencias estables antes de empezar \parencite{kruchten-rup}.
Esa premisa falla aquí dos veces: el alcance es emergente por decisión de diseño, y la capacidad disponible varía a lo largo del periodo, de modo que la red de tareas exigiría re-planificación continua, un coste permanente sin valor añadido.
A ello se suma que \ac{RUP} prescribe una disciplina de artefactos y roles pensada para equipos, desproporcionada para una persona.
De este enfoque se rescata, no obstante, una pieza: un plan macro de fases e hitos con fechas (capítulo \ref{ch:gantt-tfg}), que da la visión de conjunto que una pila de historias no ofrece.

\paragraph{Scrum puro.}
Scrum fue la metodología empleada por el autor en el proyecto grupal de la asignatura, con resultados satisfactorios, y por ello fue la primera candidata.
Sin embargo, Scrum se define explícitamente sobre un equipo con roles diferenciados (\eng{Product Owner}, \eng{Scrum Master}, \eng{Developers}) y ceremonias cuyo propósito es sincronizar y proteger a ese equipo \parencite{scrum-guide}.
En solitario, la mitad del marco queda vacío: no hay a quién hacer de contrapeso en la negociación del alcance, la \eng{daily} no tiene interlocutor y la \eng{review} no tiene público.
Además, el compromiso cerrado de cada \eng{sprint} encaja mal con una capacidad que se sabe variable: un mes de baja dedicación encadenaría «fracasos» de \eng{sprint} que no informan de nada.
Declarar «se usó Scrum» sería, en rigor, falso; lo honesto es declarar qué elementos de Scrum se adoptan y por qué.

\paragraph{Kanban puro.}
El método Kanban (flujo continuo, trabajo visualizado, límite de trabajo en curso) encaja de forma natural con la disponibilidad variable: si una semana rinde menos, fluye menos trabajo, sin ningún compromiso roto \parencite{anderson-kanban}.
Se descarta como método único porque carece de cadencia: sin iteraciones no hay ritmo fijo de re-planificación ni retrospectiva, y no se genera la medida periódica de velocidad que este \ac{TFG} necesita como telemetría para la memoria y para corregir su primera experiencia de estimación en solitario (riesgo R3).

\subsection{Decisión: método híbrido (Scrumban) adaptado a una persona}
\label{sec:metodologia-decision}

Se adopta un híbrido de la familia Scrumban \parencite{ladas-scrumban, kniberg-scrumban}: la estructura y la telemetría de Scrum que sobreviven al contexto unipersonal, sobre una mecánica de flujo Kanban que absorbe la variabilidad.
La tabla \ref{tab:metodologia-elementos} resume qué se toma de cada marco y con qué adaptación.

\begin{table}[htbp]
  \centering
  \small
  \begin{tabular}{p{3.4cm}p{1.6cm}p{8.2cm}}
    \toprule
    \textbf{Elemento} & \textbf{Origen} & \textbf{Adaptación a una persona} \\
    \midrule
    \eng{Product Backlog} & Scrum & Historias de usuario priorizadas con MoSCoW; toda idea nueva entra por aquí (control del alcance, R1). \\
    Iteraciones & Scrum & \eng{Sprints} semanales (lunes--domingo) con capacidad variable declarada por sprint. \\
    Naturaleza del compromiso & Kanban & El compromiso es el orden del \eng{backlog} y el límite de \ac{WIP}, no un contrato cerrado de \eng{sprint}; el contenido semanal es una proyección. \\
    Límite de \acs{WIP} & Kanban & Máximo 2 elementos en curso: una historia de producto y, como mucho, la épica transversal de memoria; nunca dos historias de producto a medias. \\
    Planificación de \eng{sprint} & Scrum & Sesión individual de unos 20 minutos cada lunes: fijar capacidad y seleccionar del \eng{backlog}. \\
    \eng{Daily} & Scrum & Sustituida por el procedimiento ya vigente de abrir cada sesión releyendo la última entrada del diario de desarrollo. \\
    \eng{Review} & Scrum & Verificación del incremento contra la Definición de Hecho (sección \ref{sec:dod}). \\
    Retrospectiva & Scrum & Sección fija de la entrada del diario al cierre del \eng{sprint}, que incluye la revisión de riesgos (capítulo \ref{ch:riesgos}). \\
    Velocidad & Scrum & Puntos completados por \eng{sprint}; recalibra la estimación de forma continua. \\
    \bottomrule
  \end{tabular}
  \caption{Elementos del método híbrido y su adaptación al contexto unipersonal.}
  \label{tab:metodologia-elementos}
\end{table}

La duración de una semana por \eng{sprint} (frente a las dos semanas empleadas en el proyecto grupal) se justifica por tres razones.
Primera: la realimentación de velocidad llega al doble de frecuencia justo cuando más falta hace, en las primeras semanas de la primera estimación en solitario del autor.
Segunda: el grano semanal permite modular la capacidad de septiembre sin tocar el método (cambia la capacidad, no el ritmo).
Tercera: el argumento clásico contra los \eng{sprints} cortos, el coste proporcional de las ceremonias, desaparece en solitario, donde todas las ceremonias juntas cuestan menos de una hora semanal.

\subsection{Herramienta de gestión}
\label{sec:metodologia-herramienta}

La gestión del \eng{backlog} y del tablero se realiza con GitHub Projects, con las historias como \eng{issues}, los puntos y la épica como etiquetas y las fases como \eng{milestones}.
El tablero tiene cuatro columnas (\eng{Backlog}, Preparado, En curso con límite de \acs{WIP} 2, y Hecho).
Se descarta OpenProject, la herramienta empleada en el proyecto grupal, porque su valor diferencial son las funciones de equipo (roles, actas, flujos de aprobación) que aquí no se usarían, y exige desplegar y mantener un servidor propio.
GitHub Projects, en cambio, vive junto al repositorio y enlaza cada historia con sus \eng{commits} y \eng{pull requests}, lo que genera trazabilidad historia--código sin esfuerzo adicional: evidencia gratuita para esta memoria.
El tablero es la herramienta operativa del día a día; esta memoria conserva las instantáneas relevantes (el \eng{backlog} inicial aquí, el estado final y las desviaciones en los capítulos \ref{ch:resultados} y \ref{ch:gantt-tfg}).

\section{Estimación del esfuerzo}
\label{sec:estimacion}

El esfuerzo se estima en puntos de historia (\eng{story points}), una medida relativa de tamaño, con la escala de Fibonacci modificada (1, 2, 3, 5, 8 y 13).
La estimación relativa se elige frente a la estimación directa en horas porque comparar («esto es como aquello, pero el doble») es más fiable que predecir duraciones absolutas, y porque desacopla el tamaño del trabajo de la capacidad disponible, que aquí varía por diseño \parencite{cohn-estimating}.
Al no existir equipo, no procede el \eng{planning poker}; en su lugar, cada historia se estima por triangulación contra historias de referencia ya conocidas.
La referencia inicial procede del precedente técnico del autor (el proyecto TrackBites): 2 puntos equivalen a «un \eng{endpoint} CRUD sencillo con validación compartida y pruebas de integración», y 1 punto a una tarea mecánica bien conocida.
Toda historia estimada en 13 puntos se considera demasiado grande y debe dividirse antes de entrar en un \eng{sprint} (mitigación del riesgo R3).

Como hipótesis inicial de velocidad se toman unos 11 puntos por \eng{sprint} a plena capacidad (30 horas semanales), es decir, unas 2,7 horas por punto.
Esta equivalencia sirve únicamente para el encaje macro del plan; a partir del tercer \eng{sprint} la manda la velocidad real medida, y la hipótesis se recalibra en cada retrospectiva.

\section{Definición de Hecho}
\label{sec:dod}

Una historia solo se considera terminada cuando cumple todas las condiciones siguientes:

\begin{itemize}
  \item El análisis estático y la comprobación de tipos del \eng{monorepo} pasan en verde.
  \item Las pruebas de la historia existen y pasan, y la suite completa no sufre regresiones.
  \item El cambio está integrado en la rama principal.
  \item Las decisiones tomadas quedan anotadas en el diario de desarrollo, con la alternativa descartada y su porqué.
  \item Si la historia afecta a un capítulo de esta memoria, el capítulo queda actualizado (mitigación del riesgo R4).
\end{itemize}

Esta lista es el criterio de la \eng{review} de cada incremento y protege los objetivos de calidad y trazabilidad del capítulo \ref{ch:objetivos}.

\section{\eng{Product Backlog} inicial}
\label{sec:backlog}

El \eng{backlog} inicial se construye descomponiendo las reglas de producto de \bukit en épicas, y estas en historias de usuario con el formato «como \emph{rol}, quiero \emph{acción}, para \emph{beneficio}».
Cada épica recibe una prioridad MoSCoW \parencite{clegg-moscow}, que traduce el alcance del capítulo \ref{ch:objetivos} a decisiones operativas: \emph{Must} para el alcance comprometido del que dependen los criterios de éxito, \emph{Should} para lo incluido en el alcance sin criterio de éxito asociado, \emph{Could} para los objetivos declarados secundarios, y \emph{Won't} para lo explícitamente excluido (sección \ref{sec:alcance-fuera}), que no se repite aquí.

\begin{table}[htbp]
  \centering
  \small
  \begin{tabular}{clccc}
    \toprule
    \textbf{Id} & \textbf{Épica} & \textbf{Prioridad} & \textbf{Puntos} & \textbf{Fase} \\
    \midrule
    EP01 & Fundaciones técnicas e integración continua      & Must   & 8  & F1 \\
    EP02 & Cuentas y autenticación                          & Must   & 13 & F1 \\
    EP03 & Andamiaje del cliente web y sesión               & Must   & 8  & F1 \\
    EP04 & Listas y metas                                   & Must   & 13 & F2 \\
    EP05 & Publicaciones ligadas a listas                   & Must   & 13 & F2 \\
    EP06 & Grafo social y \eng{feed} de seguidos            & Must   & 10 & F2 \\
    EP07 & Perfil y verificación de identidad               & Must   & 8  & F3 \\
    EP08 & \eng{Feed} personalizado y buscador              & Must   & 13 & F4 \\
    EP09 & Mapa social                                      & Must   & 10 & F3 \\
    EP10 & Afinidad entre personas usuarias                 & Should & 8  & --- \\
    EP11 & Mensajería y notificaciones                      & Must   & 13 & F3 \\
    EP12 & Pruebas \acs{E2E} y despliegue continuo          & Must   & 8  & F4 \\
    EP13 & \acs{IA} de estimación de plazos                 & Could  & 8  & --- \\
    EP14 & Aplicación móvil parcial                         & Could  & 13 & --- \\
    EP15 & Memoria y gestión del proyecto (transversal)     & Must   & 21 & Todas \\
    \bottomrule
  \end{tabular}
  \caption{Épicas del \eng{Product Backlog} inicial, con estimación gruesa y fase prevista.}
  \label{tab:backlog-epicas}
\end{table}

Las estimaciones de la tabla \ref{tab:backlog-epicas} son deliberadamente gruesas: siguiendo la elaboración progresiva propia de los métodos ágiles, solo se detallan en historias las épicas de la fase inminente, y el resto se refina justo antes de abordarse, con el conocimiento acumulado hasta entonces.
Como muestra del nivel de detalle al que se baja cada historia antes de entrar en un \eng{sprint}, se desarrolla una completa:

\begin{quote}
  \textbf{US-11 (EP04, 5 puntos).}
  Como persona usuaria registrada, quiero crear una lista con metas y fecha límite, para convertir mis deseos en compromisos medibles.
  \emph{Criterios de aceptación:} (1) una cuenta gratuita no puede superar las tres listas y el exceso se rechaza con un error explícito; (2) la validación vive una sola vez en el paquete compartido y la usan tanto la \ac{API} como el cliente web; (3) toda lista exige al menos una meta y una fecha límite; (4) el flujo queda cubierto por pruebas de integración.
\end{quote}

La tabla \ref{tab:backlog-f1} recoge el desglose completo de la primera fase, ya refinado.

\begin{table}[htbp]
  \centering
  \small
  \begin{tabular}{clcc}
    \toprule
    \textbf{Id} & \textbf{Historia (resumen)} & \textbf{Épica} & \textbf{Puntos} \\
    \midrule
    US-01 & Cadena de integración continua: lint, tipos y pruebas     & EP01 & 3 \\
    US-02 & Manejo centralizado de errores y registro de sucesos      & EP01 & 3 \\
    US-03 & Despliegue del esqueleto de la \acs{API}                  & EP01 & 2 \\
    US-04 & Modelo de datos de usuario y registro de cuentas          & EP02 & 3 \\
    US-05 & Inicio de sesión con \acs{JWT} de acceso y refresco       & EP02 & 3 \\
    US-06 & Rotación y revocación de \eng{refresh tokens}             & EP02 & 5 \\
    US-07 & Doble capa de limitación de peticiones                    & EP02 & 2 \\
    US-08 & Andamiaje web con enrutado de las cuatro pestañas         & EP03 & 3 \\
    US-09 & Cliente \acs{API} tipado y gestión de sesión en la web    & EP03 & 3 \\
    US-10 & Pantallas de registro e inicio de sesión                  & EP03 & 2 \\
    \bottomrule
  \end{tabular}
  \caption{Desglose en historias de la fase F1 (fundaciones), 29 puntos.}
  \label{tab:backlog-f1}
\end{table}

\section{Fases y cronograma del desarrollo}
\label{sec:fases}

El periodo de ejecución se organiza en cuatro fases que agrupan trece \eng{sprints} semanales, del 20/07/2026 al 18/10/2026, precedidas por una fase de preparación ya en curso.
Cada fase cierra con un incremento demostrable, que actúa como hito de control.

\begin{table}[htbp]
  \centering
  \small
  \begin{tabular}{clllc}
    \toprule
    \textbf{Fase} & \textbf{Contenido} & \textbf{Sprints} & \textbf{Fechas} & \textbf{Horas} \\
    \midrule
    F0 & Preparación: entorno, memoria base, este plan & ---        & hasta 19/07      & $\sim$30 \\
    F1 & Fundaciones: EP01, EP02, EP03                 & S1--S3     & 20/07 -- 09/08   & 90 \\
    F2 & Núcleo social: EP04, EP05, EP06               & S4--S6     & 10/08 -- 30/08   & 90 \\
    F3 & Diferenciadores: EP07, EP09, EP11             & S7--S11    & 31/08 -- 04/10   & 70 \\
    F4 & Cierre de producto: EP08, EP12                & S12--S13   & 05/10 -- 18/10   & 60 \\
    \bottomrule
  \end{tabular}
  \caption{Fases del desarrollo con su capacidad estimada (la épica transversal EP15 se reparte entre todas).}
  \label{tab:fases}
\end{table}

La capacidad de cada \eng{sprint} refleja la disponibilidad real esperada: 30 horas semanales en verano y octubre, y 10 horas semanales en septiembre (\eng{sprints} S7 a S10), conforme al riesgo R2 ya identificado.
Por eso la fase F3 abarca cinco \eng{sprints}: su duración en semanas no indica más trabajo, sino menos capacidad por semana.
\todonote{Confirmar las capacidades estimadas por sprint (30/10/30 horas) y la fecha de inicio de S1 antes del 20/07.}

\begin{table}[htbp]
  \centering
  \small
  \begin{tabular}{cllc}
    \toprule
    \textbf{Sprint} & \textbf{Fechas} & \textbf{Tema orientativo} & \textbf{Horas} \\
    \midrule
    S1  & 20/07 -- 26/07 & Integración continua mínima; modelo de usuario y registro & 30 \\
    S2  & 27/07 -- 02/08 & Autenticación completa con rotación de \eng{tokens}       & 30 \\
    S3  & 03/08 -- 09/08 & Andamiaje web y flujo de sesión; \emph{cierre de F1}      & 30 \\
    S4  & 10/08 -- 16/08 & Listas y metas                                            & 30 \\
    S5  & 17/08 -- 23/08 & Publicaciones desde listas                                & 30 \\
    S6  & 24/08 -- 30/08 & Grafo social y \eng{feed} de seguidos; \emph{cierre de F2} & 30 \\
    S7  & 31/08 -- 06/09 & Perfil y verificación de identidad                        & 10 \\
    S8  & 07/09 -- 13/09 & Perfil y verificación de identidad (continuación)         & 10 \\
    S9  & 14/09 -- 20/09 & Mapa social                                               & 10 \\
    S10 & 21/09 -- 27/09 & Mapa social (continuación)                                & 10 \\
    S11 & 28/09 -- 04/10 & Mensajería y notificaciones; \emph{cierre de F3}          & 30 \\
    S12 & 05/10 -- 11/10 & \eng{Feed} personalizado y buscador                       & 30 \\
    S13 & 12/10 -- 18/10 & \acs{E2E} y despliegue continuo; \emph{cierre de F4}      & 30 \\
    \bottomrule
  \end{tabular}
  \caption{Calendario de \eng{sprints} con su capacidad; el tema es una proyección, el compromiso real es el orden del \eng{backlog}.}
  \label{tab:sprints}
\end{table}

\subsection{Encaje del plan y tensión declarada}
\label{sec:fases-tension}

El inventario comprometido (\emph{Must}) suma 138 puntos, mientras que la capacidad nominal de los trece \eng{sprints} (310 horas, a la hipótesis de 2,7 horas por punto) da unos 115 puntos.
Existe por tanto una tensión inicial de unos 23 puntos ($\approx$ 60 horas, unas dos semanas a plena capacidad), que se declara aquí en lugar de maquillarse ajustando estimaciones a la baja.
La decisión consciente es triple: se mantiene el 18/10/2026 como fecha objetivo ambiciosa, se comunica una horquilla de cierre realista entre esa fecha y mediados de noviembre de 2026 (capítulo \ref{ch:gantt-tfg}), y las prioridades \emph{Should} y \emph{Could} solo entran si la velocidad real supera la hipótesis.
Esta tensión es, de hecho, la primera observación registrada del riesgo R3, y su evolución se vigila desde el primer cierre de \eng{sprint}.
En cualquier caso, la holgura estructural del calendario académico (defensa en junio de \thesisYear) hace que esta tensión amenace al objetivo propio del autor, no al \ac{TFG}.
